\section{Experimental setup}
\label{sec:setup-exp}

\paragraph{Task and data.}
Yelp sentiment transfer in the negative-to-positive direction
(\texttt{direction} $= 0\!\to\!1$). Content is measured by BERTScore against
the source; sentiment by a \texttt{bert-base-uncased} classifier fine-tuned on
Yelp.

\paragraph{Models.}
The policy is a \texttt{distilgpt2} backbone with a trained MLP adaptor head
(hidden width $2048$) that emits a fixed-length $L=5$ token prompt, with a
$\lmb$-conditioning input supplied to the head. The head is initialized with
Xavier-uniform weights at gain $10^{-4}$, which drives its output logits to
zero, so \textbf{the initial prompt distribution is numerically uniform over
the full $50{,}257$-token vocabulary}: the measured policy entropy at step $0$
is $10.824904$ against $\ln 50257 = 10.824905$. This matters for GRPO and is
revisited in Section~\ref{sec:algos-caveat}. The task model is
\texttt{gpt2-xl}, \textbf{frozen} and inference-only in \texttt{bfloat16},
served through vLLM with top-$k=10$ sampling. Only the adaptor is trained.

\paragraph{Training.}
$12{,}000$ steps; batch of $B=15$ sources per step with $\Ggrp=16$ prompts per
source ($240$ scored prompt-source rows per step, each scored with $N=50$
task-model samples); Adam at learning rate $10^{-4}$ with gradient-norm
clipping at $10$; $\lmb \sim \mathcal{U}[0,1]$ redrawn per source per step.
\rrebel{} uses $\beta = 0.5$ and refreshes its reference every $150$ steps;
GRPO uses $\beta_{\mathrm{KL}} = 0.04$ with a pinned reference. Note the two
$\beta$s are not comparable quantities --- one is a regression scale in
Eq.~\eqref{eq:logratio}, the other a KL coefficient in Eq.~\eqref{eq:grpo}.

\paragraph{Arms and seeds.}
Five arms $\times$ three seeds $=15$ runs:
three \rrebel{} variants ($\ell_1$-std, Huber-std, $\ell_1$-ent) and two GRPO
arms (plain, and with an entropy bonus $\alpha_H = 0.01$). A sixth,
\emph{superseded} GRPO configuration --- the rolling reference anchor described
in Section~\ref{sec:grpo} --- is reported separately in
Section~\ref{sec:baseline-defense} as the pre-fix baseline.

\paragraph{Evaluation protocol, and a serious caveat about it.}
Every $500$ steps the policy is evaluated over the $\lmb$ grid
$\{0, 0.1, \dots, 0.9\}$. The headline metric is \emph{mean reward}: the reward
of Eq.~\eqref{eq:reward} averaged over the evaluation set and then over the ten
$\lmb$ values; content and sentiment are averaged the same way.

The caveat is that this in-training evaluation runs on the \textbf{development
split, capped at $10$ sentences} --- not on the $500$-sentence held-out test
split. The cap is a speed hack in the data loader, and the training entry point
passes the development set as the evaluation set. \textbf{Every number in
this paper is therefore a $10$-sentence development-set number}, and no run in
this campaign was ever evaluated on the test split. Two consequences we do not
paper over: the effect sizes below carry a much wider confidence interval than
three seeds would suggest, and the absolute values are not comparable to
published test-split results. Evaluation is noisy for a second reason as well ---
the frozen task model samples stochastically (\texttt{vllm\_seed} unset during
the campaign), and re-running a single evaluation moved its score by $1.1$
points --- so we treat $\mathbf{1.0}$ \textbf{point as a noise floor} and report
differences below it as ties rather than rankings. The $6$-point family gap is
well outside that floor; the sub-$2$-point gaps between variants are not
resolvable at $n=10$.
\todo{run the $500$-sentence test-split evaluation and report it as the primary
table. \texttt{slurm/eval\_v4.slurm} does exactly this (it drives
\texttt{run\_eval.py}, which uses the test split) for the regeneration
campaign; the v3 numbers cannot be upgraded because their checkpoints are gone.}

\paragraph{Compute.}
Each run occupies one GPU ($\ge\!40$\,GB, natively bf16 --- L40S, A40, A100 or
H100) at roughly $25$\,s/step, and is executed as a chain of $4$-hour scheduler
cycles that resume from the most recent checkpoint. Generation by the frozen
task model accounts for $97.6\%$ of step time; content and sentiment scoring
together account for $2.3\%$, and the policy update for under $0.4\%$.

\paragraph{Matched-step comparison.}
Because the arms were launched at different times and restarted different
numbers of times, we compare at a matched step rather than at wall-clock
completion. Step $11{,}500$ is the largest step at which \emph{all fifteen} runs
have an evaluation, and it is the point used for the headline table; results at
step $12{,}000$ are reported alongside it with their reduced seed counts stated
explicitly (Section~\ref{sec:provenance} explains why three runs lack a
$12{,}000$ evaluation).
